\chapter*{Regression Workflow, Overfitting, and Cross-Validation}
\addcontentsline{toc}{chapter}{Regression Workflow, Overfitting, and Cross-Validation}

\section*{A Regression Example: End-to-End Workflow}

To ground the abstract ERM framework in a concrete workflow, consider a canonical regression task. We observe $m = 10$ noisy samples from an unknown function $f(x) = \sin(2\pi x)$, where $x \in [0,1]$ and $t \in [-1,1]$. The dataset is:
$$D = \{(x_i, t_i)\}_{i=1}^{m}, \quad t_i = \sin(2\pi x_i) + \epsilon_i$$
where $\epsilon_i$ is random noise. The goal: learn $f: X \to T$ from $D$ alone---without knowing the true generating function.

\subsection*{Step 1: Choose the Hypothesis Class $\mathcal{H}$}

We hypothesize that $f$ belongs to the family of \textbf{polynomial functions of degree $M$}:
$$f_{\boldsymbol{\theta}}(x) = \theta_0 + \theta_1 x + \theta_2 x^2 + \cdots + \theta_M x^M = \sum_{j=0}^{M} \theta_j x^j$$
with parameter vector $\boldsymbol{\theta} \in \mathbb{R}^{M+1}$. Two critical observations:
\begin{itemize}
    \item $M$ controls the \textbf{complexity/capacity} of the hypothesis class. Higher $M$ means more expressive models.
    \item The model is \textbf{linear in its parameters} $\boldsymbol{\theta}$ (even though it is non-linear in $x$), so standard linear algebra solvers apply.
    \item A degree-$M$ polynomial has $M+1$ free parameters to learn.
\end{itemize}

\subsection*{Step 2: Define the Loss Function}

\begin{definition}[Sum of Squared Errors (SSE)]
$$\text{SSE}(\boldsymbol{\theta}) = \frac{1}{2} \sum_{i=1}^{m} \left( f_{\boldsymbol{\theta}}(x_i) - t_i \right)^2$$
\end{definition}

\begin{definition}[Root Mean Square Error (RMSE)]
$$\text{RMSE}(\boldsymbol{\theta}) = \sqrt{\frac{\text{SSE}(\boldsymbol{\theta})}{m}}$$
RMSE normalizes by dataset size and has the same units as the target, making it interpretable and comparable across datasets of different sizes. This serves as our empirical loss $L_{\text{emp}}$.
\end{definition}

\subsection*{Step 3: Train (Minimize $L_{\text{emp}}$)}

For each choice of $M$, we solve $\boldsymbol{\theta}^* = \arg\min_{\boldsymbol{\theta}} \text{SSE}(\boldsymbol{\theta})$ using least-squares. The following snippet fits polynomials of varying degrees to the $\sin(2\pi x)$ dataset:

\begin{lstlisting}[style=pythonstyle]
import numpy as np
from numpy.polynomial import polynomial as P

# Generate noisy sin(2*pi*x) data
np.random.seed(0)
x = np.linspace(0, 1, 10)
t = np.sin(2 * np.pi * x) + np.random.randn(10) * 0.3

# Fit polynomials of degree M = 0, 1, 3, 9
for M in [0, 1, 3, 9]:
    coeffs = np.polyfit(x, t, M)    # Least-squares fit
    p = np.poly1d(coeffs)           # Callable polynomial
    residuals = p(x) - t
    rmse = np.sqrt(np.mean(residuals**2))
    print(f"M={M}: RMSE={rmse:.4f}, params={M+1}")
\end{lstlisting}

\begin{remark}[Parameters vs.\ Data Points]
A degree-$M$ polynomial has $M+1$ parameters. With $m = 10$ data points: $M = 0$ has 1 parameter (constant), $M = 1$ has 2 (line), $M = 3$ has 4 (cubic), and $M = 9$ has 10---exactly as many parameters as data points.
\end{remark}

\section*{Overfitting vs.\ Generalization}

Training error ($L_{\text{emp}}$ on $D$) tells us how well the model fits the data it has seen. But the real question is: \textit{how will it perform on new, unseen data?}

\subsection*{The Core Phenomenon}

As model complexity $M$ increases:
\begin{enumerate}
    \item \textbf{Training RMSE decreases monotonically.} More parameters $\Rightarrow$ more capacity to fit the training points. At $M = 9$ with $m = 10$ points, the polynomial passes through \textit{every} data point: $\text{RMSE}_{\text{train}} = 0$.
    \item \textbf{Validation RMSE first decreases, then rapidly increases.} The model begins fitting the noise rather than the underlying signal.
\end{enumerate}

\begin{definition}[Overfitting]
\textbf{Overfitting} occurs when a model fits too closely to the training data---learning the noise, not the signal---resulting in excellent training performance but poor generalization to unseen data. Overfitting defeats the fundamental purpose of ML.
\end{definition}

\begin{remark}[When Does Overfitting Occur?]
Overfitting depends on the relationship between \textbf{model complexity} and \textbf{data quantity}:
\begin{itemize}
    \item When $M \geq m$ (parameters $\geq$ data points), the model can interpolate every training point perfectly, almost certainly overfitting.
    \item When $N \gg M$ (much more data than parameters), even complex models can generalize well---there is enough signal to constrain the parameters.
    \item Simpler models tend to prevent overfitting and generalize better. This principle motivates \textbf{regularization} and connects to Occam's razor: prefer the simplest hypothesis consistent with the data.
\end{itemize}
\end{remark}

\subsection*{The Sweet Spot}

The optimal model complexity corresponds to the point where:
\begin{itemize}
    \item Validation error is near its \textbf{minimum}, and
    \item Training error is reasonably low (but \textit{not} zero---a zero training error on noisy data is a red flag).
\end{itemize}

\section*{Model Selection: The Bias-Variance Tradeoff}

The tension between underfitting and overfitting is formalized by the \textbf{bias-variance tradeoff}:

\begin{description}
    \item[Bias:] Error from overly simplistic assumptions in the hypothesis class. High bias $\Rightarrow$ the model systematically misses the true pattern (\textbf{underfitting}).
    \item[Variance:] Error from excessive sensitivity to the specific training set. High variance $\Rightarrow$ the model changes drastically with different samples (\textbf{overfitting}).
\end{description}

Increasing complexity \textit{reduces} bias but \textit{increases} variance. The goal of model selection is to find the complexity that minimizes the \textbf{total} generalization error, which is approximately $\text{Bias}^2 + \text{Variance} + \text{irreducible noise}$.

\section*{The Hold-Out Method}

The simplest approach to estimating generalization error: split the available data into two disjoint sets.

\begin{definition}[Hold-Out Method]
Partition the dataset into:
\begin{itemize}
    \item \textbf{Training set} $D$: used to learn parameters $\boldsymbol{\theta}^*$.
    \item \textbf{Test set} $T$: used to evaluate $L_{\text{emp}}(T)$ as an estimate of $L_{\text{gen}}$.
\end{itemize}
The fundamental rationale: evaluate the model on data \textbf{independent} from the data used for training.
\end{definition}

\begin{remark}[Drawbacks of Hold-Out]
\begin{itemize}
    \item \textbf{Wastes data:} test examples are never used for training, which hurts performance when data is scarce.
    \item \textbf{High variance:} a single (un)fortunate split can produce a misleading estimate of generalization error.
    \item \textbf{Split dependency:} results can vary significantly depending on which examples end up in which set.
\end{itemize}
\end{remark}

\section*{Cross-Validation (CV)}

\textbf{Cross-validation} overcomes the limitations of hold-out by iterating the split-train-evaluate procedure $K$ times, rotating which data plays the training and validation roles.

\subsection*{$K$-Fold Cross-Validation}

\begin{definition}[$K$-Fold CV]
Partition the dataset into $K$ equal-sized \textbf{folds} $F_1, F_2, \ldots, F_K$. For each fold $k = 1, \ldots, K$:
\begin{enumerate}
    \item \textbf{Train} on all folds except $F_k$: $D_{\text{train}} = D \setminus F_k$.
    \item \textbf{Validate} on $F_k$: compute $L_{\text{emp}}(F_k)$.
\end{enumerate}
The CV estimate of generalization error is the average across folds:
$$L_{\text{CV}} = \frac{1}{K} \sum_{k=1}^{K} L_{\text{emp}}(F_k)$$
\end{definition}

Every data point serves as a validation example exactly once and as a training example $K-1$ times. Common choice: $K = 5$ or $K = 10$.

\subsection*{Variants}

\begin{description}
    \item[Leave-One-Out CV (LOOCV):] The extreme case $K = N$. Each fold contains a single example. Maximizes training data per fold but is computationally expensive ($N$ training runs) and can have high variance for small datasets.
    \item[Random Subsampling CV:] Instead of deterministic partitions, randomly sample training/validation splits each iteration. More flexible but some points may never appear in validation.
\end{description}

\begin{remark}[Formal Guarantees]
Under mild conditions and reasonably large datasets, CV provides not just an empirical estimate of generalization error, but comes with formal bounds relating the CV estimate to the true (unknown) $L_{\text{gen}}$.
\end{remark}

\subsection*{$K$-Fold CV with \texttt{sklearn}}

\begin{lstlisting}[style=pythonstyle]
import numpy as np
from sklearn.preprocessing import PolynomialFeatures
from sklearn.linear_model import LinearRegression
from sklearn.pipeline import make_pipeline
from sklearn.model_selection import cross_val_score

# Generate noisy sin(2*pi*x) data
np.random.seed(0)
X = np.linspace(0, 1, 30).reshape(-1, 1)
y = np.sin(2 * np.pi * X).ravel() + np.random.randn(30) * 0.3

# Compare polynomial degrees via 5-fold CV
for M in [1, 3, 5, 9]:
    model = make_pipeline(
        PolynomialFeatures(M), LinearRegression()
    )
    # neg_mean_squared_error: sklearn negates so higher=better
    scores = cross_val_score(
        model, X, y, cv=5, scoring='neg_mean_squared_error'
    )
    rmse_scores = np.sqrt(-scores)
    print(f"M={M}: CV RMSE = {rmse_scores.mean():.3f}"
          f" (+/- {rmse_scores.std():.3f})")
\end{lstlisting}

\section*{The General Model Selection Pipeline}

In practice, model selection combines hold-out and cross-validation in a principled pipeline:

\begin{enumerate}
    \item \textbf{Initial split (Hold-Out):} Partition the full dataset into a \textbf{Training set} and a \textbf{Test set}. The test set is locked away---untouched until the very end.
    \item \textbf{Internal CV:} Further split the Training set into internal training folds and validation folds using $K$-Fold CV.
    \item \textbf{Train and validate candidates:} For each candidate model (varying hyperparameters, algorithms, or complexity), run the full CV loop. Record each model's estimated generalization error $L_{\text{CV}}$.
    \item \textbf{Select the best model:} Choose the candidate with the lowest $L_{\text{CV}}$.
    \item \textbf{Retrain on full Training set:} Retrain the selected model using the \textit{entire} Training set (all folds combined) to maximize the data available for learning.
    \item \textbf{Final evaluation on Test set:} Evaluate the retrained model on the held-out Test set. This is the unbiased estimate of deployment performance. If the test error is much larger than $L_{\text{CV}}$, reconsider the candidate models or data pipeline.
\end{enumerate}

\begin{remark}[Why Retrain?]
After CV selects the best model \textit{type} (e.g., $M = 3$), we retrain on all training data because CV only used $\frac{K-1}{K}$ of it per fold. The final model benefits from every available training example.
\end{remark}

\begin{remark}[Test Set Discipline]
The test set must remain untouched during model selection. If you use test performance to choose between models, it becomes a validation set---and your generalization estimate becomes optimistically biased. This is a common and serious methodological error.
\end{remark}
